\chapter{GIỚI THIỆU ĐỀ TÀI}

\section{Đặt vấn đề}

Các mô hình ngôn ngữ lớn hiện nay không chỉ được sử dụng để sinh văn bản mà còn có thể tham gia các vòng lặp hành động để lập kế hoạch, gọi công cụ và tương tác với dịch vụ bên ngoài~\cite{Yao2022ReAct,Schick2023Toolformer}. Khi được triển khai dưới dạng tác nhân AI, phần mềm có thể tự động chia nhiệm vụ thành nhiều bước, gọi API nhiều lần, đánh giá kết quả trung gian và thử lại khi gặp lỗi. Mỗi bước như vậy có thể phát sinh chi phí sử dụng dịch vụ.

Điểm khác biệt của tác nhân AI so với người dùng thông thường nằm ở tính tự động và tốc độ thực thi. Một người dùng thường chủ động xác nhận từng thao tác có chi phí, trong khi một tác nhân AI có thể tạo nhiều yêu cầu liên tiếp hoặc đồng thời mà không có xác nhận thủ công ở từng bước. Nếu không có cơ chế kiểm soát ở giữa, một vòng lặp sai, một chiến lược thử lại không phù hợp hoặc nhiều lời gọi song song có thể nhanh chóng làm vượt ngân sách được cấp.

Các nhà cung cấp dịch vụ AI thường tính phí theo mô hình, số token hoặc loại tài nguyên đã sử dụng. Chi phí thực tế chỉ được xác định sau khi yêu cầu hoàn tất. Nếu hệ thống chỉ ghi nhận chi phí sau khi nhận kết quả, nhiều yêu cầu đồng thời có thể cùng vượt qua bước kiểm tra số dư và làm phát sinh chi tiêu vượt mức. Việc lựa chọn mô hình theo chi phí có thể giảm tổng chi phí sử dụng~\cite{Chen2023FrugalGPT}, nhưng chưa giải quyết đầy đủ nhu cầu kiểm soát ngân sách trước thực thi và truy vết hành vi của từng tác nhân AI.

Từ thực tế đó, đề tài tập trung thiết kế và xây dựng Asymptotic, một \textit{AI Agent Financial Gateway} đóng vai trò trung gian giữa tác nhân AI bên ngoài và các nhà cung cấp dịch vụ AI trả phí. Tổ chức duy trì ví và phân bổ ngân sách theo chuỗi tổ chức, đội ngũ, lập trình viên và tác nhân AI. Lập trình viên đăng ký tác nhân do mình phát triển hoặc quản lý, cấp hạn mức trong phạm vi được nhận và quản lý khóa truy cập của tác nhân. Khi sử dụng dịch vụ AI, tác nhân gọi Gateway bằng khóa do Asymptotic cấp; Gateway kiểm tra định danh, chính sách và ngân sách trước khi gọi nhà cung cấp bằng thông tin xác thực nội bộ.

Trong đề tài, kiểm soát theo thời gian thực được hiểu là việc Gateway đưa ra quyết định chấp nhận hoặc từ chối đồng bộ trên từng yêu cầu trước khi yêu cầu đó phát sinh chi phí tại nhà cung cấp. Cách tiếp cận này khác với việc chỉ tổng hợp hóa đơn hoặc cảnh báo sau khi dịch vụ đã được sử dụng.

\section{Bài toán cần giải quyết}

Đề tài cần giải quyết bài toán kiểm soát chi phí tại thời điểm tác nhân AI phát sinh yêu cầu. Trước khi gọi nhà cung cấp, hệ thống phải xác định tác nhân tạo yêu cầu, lập trình viên quản lý, đội ngũ và tổ chức chịu chi phí; đồng thời kiểm tra khóa truy cập, chính sách, hạn ngạch và hạn mức khả dụng trên toàn bộ đường dẫn ngân sách.

Các vấn đề chính bao gồm:

\begin{itemize}
    \item \textbf{Định danh tác nhân AI:} Mỗi tác nhân bên ngoài cần được lập trình viên có quyền đăng ký vào Gateway và gắn với tổ chức, đội ngũ cùng người quản lý cụ thể.
    \item \textbf{Quản lý quyền truy cập:} Khóa do Asymptotic cấp cho tác nhân phải được quản lý tách biệt với thông tin đăng ký tác nhân và với thông tin xác thực nội bộ dùng để gọi nhà cung cấp.
    \item \textbf{Quản lý ngân sách:} Tổ chức duy trì tiền trong ví; quản trị viên tổ chức cấp hạn mức cho đội ngũ và lập trình viên; lập trình viên cấp hạn mức cho các tác nhân do mình quản lý.
    \item \textbf{Kiểm soát trước khi phát sinh chi phí:} Yêu cầu chỉ được chuyển tiếp tới nhà cung cấp khi thỏa mãn chính sách sử dụng và ngân sách khả dụng.
    \item \textbf{Định tuyến an toàn tới nhà cung cấp:} Tác nhân không dùng trực tiếp khóa của nhà cung cấp AI; Gateway gọi nhà cung cấp bằng thông tin xác thực nội bộ do hệ thống quản lý.
    \item \textbf{Truy vết hành vi và chi phí:} Mỗi lần sử dụng dịch vụ cần được liên kết với tổ chức, tác nhân, nhà cung cấp, mô hình, chi phí và trạng thái xử lý.
\end{itemize}

Khó khăn cốt lõi nằm ở việc chi phí cuối cùng chỉ được biết sau khi nhà cung cấp xử lý yêu cầu, trong khi nhiều yêu cầu có thể đồng thời sử dụng cùng một nguồn ngân sách. Hệ thống phải kiểm soát hạn mức nhất quán, hạn chế chi tiêu vượt mức, tránh ghi nhận trùng khi yêu cầu được gửi lại và lưu đủ dữ liệu để đối soát sau khi xử lý.

\section{Mục tiêu và phạm vi}

Mục tiêu của đề tài là thiết kế và xây dựng nguyên mẫu \textit{AI Agent Financial Gateway} để kiểm soát quyền truy cập, ngân sách và chi phí của các yêu cầu sử dụng dịch vụ AI có tính phí do tác nhân AI bên ngoài phát sinh. Nguyên mẫu hướng tới các mục tiêu sau:

\begin{itemize}
    \item Cung cấp một điểm vào thống nhất để tác nhân AI bên ngoài gọi dịch vụ AI trả phí thông qua Asymptotic;
    \item Cho phép lập trình viên đăng ký, quản lý tác nhân và quản lý khóa truy cập do Asymptotic cấp cho từng tác nhân;
    \item Hỗ trợ nạp tiền vào ví tổ chức và phân bổ hạn mức theo chuỗi tổ chức, đội ngũ, lập trình viên và tác nhân AI;
    \item Kiểm tra khóa, chính sách, hạn ngạch và ngân sách trước khi định tuyến yêu cầu hợp lệ tới nhà cung cấp bằng thông tin xác thực nội bộ;
    \item Hạn chế tính phí trùng khi yêu cầu được gửi lại và ghi nhận mức sử dụng, chi phí, giao dịch cùng dấu vết thực thi để phục vụ kiểm tra, báo cáo và đối soát;
    \item Đo riêng thời gian xử lý phát sinh tại Gateway để đánh giá chi phí độ trễ của lớp kiểm soát.
\end{itemize}

Phạm vi MVP bao gồm giao diện lập trình phía máy chủ; quản lý tổ chức, đội ngũ và lập trình viên; ví tổ chức và chuỗi hạn mức; đăng ký tác nhân AI bên ngoài; quản lý khóa truy cập; cấu hình nhà cung cấp, mô hình và biểu giá; kiểm tra chính sách và ngân sách trước thực thi; chuyển tiếp yêu cầu; kiểm soát lũy đẳng; ghi nhận mức sử dụng, chi phí, giao dịch và dấu vết thực thi. Việc nạp tiền được triển khai hoặc mô hình hóa ở mức đủ để chứng minh luồng tài chính cốt lõi.

Hệ thống không tạo, huấn luyện, triển khai, lưu trú hoặc điều phối tác nhân AI. Tác nhân là phần mềm bên ngoài do lập trình viên phát triển hoặc quản lý, sau đó được lập trình viên có quyền đăng ký vào Gateway. Quản lý hạn mức của tác nhân và quản lý khóa truy cập là các chức năng riêng, không nằm trong thao tác đăng ký tác nhân.

Các nội dung ngoài phạm vi MVP gồm cổng thanh toán hoàn chỉnh, hóa đơn và quyết toán ở mức vận hành thực tế, bảng điều khiển quản trị đầy đủ, tối ưu định tuyến nhiều mô hình theo thời gian thực, điều phối đa tác nhân và triển khai phân tán có tính sẵn sàng cao.

\section{Phương pháp thực hiện}

Đề tài được thực hiện qua sáu bước: (1) khảo sát đặc điểm thực thi của tác nhân AI và rủi ro chi phí khi sử dụng dịch vụ AI trả phí; (2) xác định actor, yêu cầu chức năng, yêu cầu phi chức năng và các trường hợp sử dụng; (3) đặc tả trường hợp sử dụng và mô hình hóa luồng nghiệp vụ; (4) phân tích lớp theo mô hình Boundary, Control và Entity; (5) thiết kế kiến trúc, gói, lớp, tương tác, trạng thái, dữ liệu và giao diện API; (6) triển khai, kiểm thử nguyên mẫu và đánh giá mức độ đáp ứng yêu cầu.

\section{Bố cục đồ án}

Đồ án gồm sáu chương:

Chương 1 giới thiệu vấn đề, bài toán, mục tiêu, phạm vi và phương pháp thực hiện. Chương 2 khảo sát hiện trạng, đặc tả yêu cầu, phân tích actor, phân tích và đặc tả các trường hợp sử dụng, đồng thời mô hình hóa luồng nghiệp vụ chính. Chương 3 phân tích hướng đối tượng, nhận diện và đối chiếu các lớp Boundary, Control, Entity với trường hợp sử dụng.

Chương 4 trình bày kiến trúc và thiết kế hướng đối tượng, bao gồm gói, lớp, tương tác, trạng thái, dữ liệu và giao diện API. Chương 5 mô tả môi trường triển khai, chức năng đã triển khai, kiểm thử, đánh giá kết quả và hạn chế triển khai. Chương 6 tổng kết kết quả, hạn chế và hướng phát triển của đồ án.

\section*{Kết luận chương}

Chương này đã xác định nhu cầu xây dựng một cổng tài chính thời gian thực dành riêng cho các yêu cầu do tác nhân AI tự động phát sinh. Trọng tâm của đề tài là kiểm soát định danh, hạn mức, ngân sách, chi phí và lịch sử sử dụng của tác nhân AI khi tương tác với các dịch vụ AI có tính phí. Các mục tiêu và phạm vi nêu trên là cơ sở để phân tích yêu cầu hệ thống trong Chương 2.
